% SEGMENT OLP-0010-S01
% Part: sets-functions-relations
% Chapter: sets
% Section: russells-paradox

\documentclass[../../../include/open-logic-section]{subfiles}

\begin{document}

\olfileid{sfr}{set}{rus}
\olsection{ரஸ்ஸலின் முரண்பாடு}

$\phi(x)$ என்ற நிபந்தனையை நிறைவேற்றும் $x$ களின் \emph{அந்தக்} கணத்தைக்
குறிக்க $\Setabs{x}{\phi(x)}$ என்ற குறியீட்டைப் பயன்படுத்த
உறுப்புசார் சமத்துவம் வழிவகுக்கிறது. எனினும், அது \emph{உண்மையில்}
நியாயப்படுத்துவது பின்வரும் கருத்தை மட்டுமே.
$\phi$ என்ற பண்புடைய அனைத்தையும், அவற்றை மட்டுமே, உறுப்புகளாகக் கொண்ட
ஒரு கணம் \emph{இருந்தால்}, \emph{அப்போது} அத்தகைய கணம் ஒன்றே ஒன்றாக இருக்கும்.
வேறுவிதமாகச் சொன்னால், ஒரு $\phi$ ஐ நிர்ணயித்த பிறகு,
$\Setabs{x}{\phi(x)}$ என்னும் கணம் \emph{இருக்குமானால்}, அது தனித்ததாகும்.

% SEGMENT OLP-0010-S02
ஆனால் இந்த நிபந்தனை முக்கியமானது! எல்லாப் பண்புகளும்
\emph{பண்புவழிக் கணமாக்கலுக்கு} இடமளிப்பதில்லை.
அதாவது, சில பண்புகள் கணங்களை வரையறுப்பது \emph{இல்லை}.
அனைத்துப் பண்புகளும் அவ்வாறு வரையறுத்தால், நேரடியான முரண்கள் ஏற்படும்.
இதன் மிகவும் புகழ்பெற்ற எடுத்துக்காட்டு ரஸ்ஸலின் முரண்பாடாகும்.

% SEGMENT OLP-0010-S03
கணங்கள் பிற கணங்களின் !!{element}s ஆக இருக்கலாம்;
எடுத்துக்காட்டாக, $A$ என்னும் கணத்தின் அடுக்குக் கணம் கணங்களால் ஆனது.
ஆகவே, ஒரு கணம் மற்றொரு கணத்தின் !!a{element} ஆக இருக்கிறதா எனக் கேட்பதற்கும்
ஆராய்வதற்கும் பொருள் உண்டு. ஒரு கணம் தனக்குத்தானே உறுப்பாக இருக்க முடியுமா?
கணம் என்ற கருத்தில் எதுவும் இதை விலக்குவதாகத் தெரியவில்லை.
எடுத்துக்காட்டாக, \emph{அனைத்துக்} கணங்களும் பொருள்களின் ஒரு தொகுப்பாக
அமைந்தால், அவற்றை ஒரே கணமாக, அதாவது அனைத்துக் கணங்களின் கணமாக,
ஒன்றுதிரட்டலாம் என ஒருவர் நினைக்கக்கூடும். அதுவும் ஒரு கணமாக இருப்பதால்,
அனைத்துக் கணங்களின் கணத்தில் அதுவே !!a{element} ஆக இருக்கும்.

% SEGMENT OLP-0010-S04
தன்னைத்தானே !!a{element} ஆகக் கொண்டிராத பண்பை, அதாவது
\emph{தன்னுறுப்பில்லாததாக} இருக்கும் பண்பைக் கருதும்போது
ரஸ்ஸலின் முரண்பாடு எழுகிறது. தம்மைத்தாமே !!a{element} ஆகக் கொண்டிராத
அனைத்துக் கணங்களையும் கொண்ட ஒரு கணம் இருப்பதாகக் கொண்டால் என்ன ஆகும்?
\[
R = \Setabs{x}{x \notin x}
\]
என்பது இருக்கிறதா? அது இல்லை என்பதை நிறுவ முடியும்.

% SEGMENT OLP-0010-S05
\begin{thm}[ரஸ்ஸலின் முரண்பாடு]\ollabel{thm:russells-paradox}
	$R = \Setabs{x}{x \notin x}$ எனும் கணம் எதுவும் இல்லை.
\end{thm}

% SEGMENT OLP-0010-S06
\begin{proof}
$R = \Setabs{x}{x \notin x}$ இருந்தால்,
$R \notin R$ ஆக இருக்கும்போது, அப்போதும் அப்போது மட்டுமே $R \in R$.
இது ஒரு முரணாகும்.
\end{proof}

% SEGMENT OLP-0010-S07
\begin{tagblock}{novice}
\begin{explain}
இந்த நிறுவலை மேலும் மெதுவாகப் பார்ப்போம். $R$ இருந்தால்,
$R \in R$ ஆகிறதா இல்லையா எனக் கேட்பதற்குப் பொருள் உண்டு.
$R \in R$ என்பதே உண்மை எனக் கொள்வோம்.
தம்மைத்தாமே !!{element}s ஆகக் கொண்டிராத அனைத்துக் கணங்களின் கணமாக
$R$ வரையறுக்கப்பட்டது. எனவே, $R \in R$ எனில்,
$R$ ஐ வரையறுக்கும் பண்பு $R$ க்கே இல்லை.
ஆனால் அந்தப் பண்பைக் கொண்ட கணங்கள் மட்டுமே $R$ இல் உள்ளன.
ஆகவே $R$ என்பது $R$ இன் !!a{element} ஆக இருக்க முடியாது;
அதாவது $R \notin R$. ஆனால் $R$ என்பது $R$ இன் !!a{element} ஆகவும்
அவ்வாறு இல்லாமலும் ஒரே நேரத்தில் இருக்க முடியாது.
எனவே ஒரு முரண் கிடைக்கிறது.

% SEGMENT OLP-0010-S08
$R \in R$ என்ற எடுகோள் முரணுக்கு இட்டுச் செல்வதால், $R \notin R$.
ஆனால் இதுவும் முரணுக்கு இட்டுச் செல்கிறது!
ஏனெனில் $R \notin R$ எனில், $R$ ஐ வரையறுக்கும் பண்பு $R$ க்கே உள்ளது.
எனவே தன்னுறுப்பில்லாத பிற கணங்களைப் போலவே, $R$ என்பதும் $R$ இன்
!!a{element} ஆக இருக்கும். மீண்டும், அது $R$ இன் !!a{element} ஆக இல்லாமலும்
அவ்வாறு இருக்கவும் ஒரே நேரத்தில் முடியாது.
\end{explain}
\end{tagblock}

% SEGMENT OLP-0010-S09
\begin{digress}
ரஸ்ஸலின் முரண்பாட்டில் சிக்காத ஒரு கணக்கோட்பாட்டை எவ்வாறு அமைப்பது?
அதாவது, $R = \Setabs{x}{x \notin x}$ இருக்கிறது என்ற
\emph{முரணுடைய} கூற்றைத் தவிர்ப்பது எப்படி?
கணங்கள் எப்போது இருக்கின்றன, எப்போது இருப்பதில்லை என்பதைச் சொல்லும்
மிகத் துல்லியமான நிபந்தனைகளைத் தரும் அடிக்கோள்களை வகுக்க வேண்டும்.

% SEGMENT OLP-0010-S10
இந்த அத்தியாயத்தில் சுருக்கமாக விவரிக்கப்பட்ட கணக்கோட்பாடு அவ்வாறு செய்யவில்லை.
அது \emph{உண்மையிலேயே முறைசாராதது}.
கணங்கள் உறுப்புசார் சமத்துவத்தை நிறைவேற்றுகின்றன என்பதையும்,
சில கணங்கள் தரப்பட்டால் அவற்றின் சேர்ப்பு, வெட்டு முதலியவற்றை
அமைக்கலாம் என்பதையும் மட்டுமே அது கூறுகிறது.
இதைக் காட்டிலும் அதிகத் துல்லியத்துடன் கணக்கோட்பாட்டை உருவாக்க முடியும்.
\oliflabeldef{cumul:::part}{அந்தத் துல்லியமான அணுகுமுறையை
\olref[cumul][][]{part} என்னும் பகுதிக்கு ஒதுக்குவோம்.
இப்போதைக்கு முறைசாராமல் தொடர்ந்து, முரண்களை கவனமாகத் தவிர்க்க முயல்வோம்.}{}
\end{digress}

\end{document}
